\chapter{SCHEME MODEL}
\graphicspath{{Chapter2/}}

\section{System Architecture}
The proposed asymmetric searchable encryption scheme includes mainly three entities, namely, the Data Owner (DO), the Cloud Server (CS), and the Users. Unlike the traditional symmetric-based schemes, which often involve complex key distribution issues, the proposed scheme exploits the power of asymmetric cryptography. 

\subsection{System Entities}
\begin{itemize}
    \item \textbf{Data Owner (DO):} The DO is responsible for generating the system's cryptographic keys, specifically a Public Key (PK) and a Master Secret Key (MSK). The data owner utilizes the PK to generate secure Bloom filter indexes and encrypt the files prior to uploading them to the cloud environment. 
    Crucially, the DO acts as the central authority that uses the MSK to grant search capabilities to authorized users by generating secure search tokens.
    \item \textbf{Cloud Server (CS):} The CS is an entity with abundant storage and computational resources. It serves as the repository for the ciphertexts and their corresponding searchable indexes. Upon taking in a user's search token, the server runs the IPE evaluation routine against the index structures to identify and retrieve the appropriate documents.
    \item \textbf{Users:} Users are entities authorized by the DO to perform keyword searches over the outsourced encrypted data. To trigger a retrieval request, an authorized individual maps their search terms to a query vector, obtains a specific search token from the DO, and forwards this token to the storage server.
\end{itemize}

\subsection{Operational Flow}
The workflow of the proposed scheme operates in the following phases:
\begin{enumerate}
    \item \textbf{Setup:} The DO initializes the IPE system parameters, publishing the PK and securing the MSK.
    \item \textbf{Data Outsourcing:} The DO isolates terms from every file, maps them to bigram arrays, and populates a Bloom filter utilizing LSH functions \cite{wang2014privacy}.% [cite: 139, 152] 
    This filter is then secured with the PK, creating the final index that accompanies the encrypted file to the CS.% [cite: 62]
    \item \textbf{Trapdoor Generation:} An authorized user wishes to search for a set of keywords. The keywords are transformed into a query vector and sent to the DO. The DO applies the MSK to generate a secure search token ($TK_{\mathcal{Q}}$) and returns it to the user.
    \item \textbf{Secure Search via MP-LSH:} The user submits $TK_{\mathcal{Q}}$ to the CS. The CS applies the IPE evaluation algorithm \cite{shen2009predicate} to compute the inner product between the token and the stored encrypted indexes. To resolve the high false negative rate of standard LSH, the CS utilizes Multi-Probe LSH (MP-LSH) logic \cite{lv2007multi}, systematically checking adjacent hash buckets to locate fuzzy matches that fall slightly outside the strict hashing boundaries. The CS then returns the matching encrypted files.% [cite: 64]
\end{enumerate}

\section{Security Model}
We consider the server to be "honest but curious" with respect to the service provider.% [cite: 66] 
With these assumptions, the server reliably adheres to its operational protocols and provides data availability.% [cite: 67] 
However, the server plays the role of a passive eavesdropper, which might allow it to examine the ciphertexts, indexes, and queries to gain insights into the confidential file information or the behavior of users searching for the files.% [cite: 68, 78]

In order to thoroughly test the privacy of the system, we consider two types of adversaries based on the level of intelligence they possess:% [cite: 79]
\begin{itemize}
    \item \textbf{Known Ciphertext Model:} The adversary is restricted to viewing the outsourced ciphertexts, the encrypted index vectors, and the submitted search tokens.% [cite: 80] 
    The server tracks access patterns---noting which tokens yield which files---but remains ignorant of any outside contextual information.% [cite: 81]
    \item \textbf{Known Background Model:} The adversary is equipped with supplementary statistical knowledge about the outsourced collection, including term frequencies or historical distributions derived from comparable plaintext archives.% [cite: 83, 84, 85] 
    The adversary attempts to use this statistical knowledge to launch linear analysis attacks against the secure indexes to deduce the underlying keywords.% [cite: 288, 290]
\end{itemize}

\section{Design Goals}
To successfully overcome the hurdles of encrypted search in shared ecosystems, our asymmetric MP-LSH architecture pursues a distinct set of operational and security targets:% [cite: 102]
\begin{itemize}
    \item \textbf{Multi-Keyword Fuzzy Search:} The architecture must handle complex, multi-term queries and gracefully manage typographical mistakes, guaranteeing that relevant files are surfaced despite minor spelling anomalies.% [cite: 103, 104]
    \item \textbf{High Result Accuracy (Recall):} Through the adoption of MP-LSH \cite{lv2007multi}, our design seeks to dramatically lower the instances of missed matches inherent to basic LSH systems, maximizing retrieval recall.% [cite: 106, 107]
    \item \textbf{Robust Privacy Guarantees:} Our framework guarantees the protection of sensitive file contents and search terms.% [cite: 105] 
    The storage provider remains incapable of extracting plaintext terms from the stored indexes or queries (Keyword Privacy), and it is prevented from linking distinct tokens that represent identical search terms (Trapdoor Unlinkability).% [cite: 91, 93]
    \item \textbf{Dynamic Updates Without Dictionaries:} Moving away from legacy systems that depend on rigid, global vocabulary lists, our architecture enables fluid data modifications---such as insertions, deletions, or edits---without requiring a full reconstruction of the searchable index.% [cite: 108, 109, 219]
\end{itemize}